\section{Force Constants}
\label{sec:fd}
% ----------------------------------------------------------------------

\subsection{Potential-energy expansion}

For theoretical description of lattice vibrations,  the crystal potential energy
in the atomic displacements is expanded from equilibrium
\cite{born_dynamical_1996, maradudin_theory_1963},
\begin{multline}\label{eq:e_exp_ph}
    E\bigl(\{u_i\}_{i=1}^{\ndof}\bigr)
    = E_0 +
      \sum_{i}^{\ndof} \frac{\partial E}{\partial u_i}\, u_i +
      \frac{1}{2}\sum_{ij}^{\ndof}
        \frac{\partial^2 E}{\partial u_i\,\partial u_j}\, u_i\, u_j \\
    + \frac{1}{6}\sum_{ijk}^{\ndof}
        \frac{\partial^3 E}{\partial u_i\,\partial u_j\,\partial u_k}\,
        u_i\, u_j\, u_k
    + \cdots,
\end{multline}
where $\ndof = 3\nsuper = 3\Ncell\nprim$ is the total number of
degrees of freedom (see \cref{sec:notation_fourier} for the notation).
At equilibrium, the first-order term vanishes.  The
second-order coefficients define the harmonic, or second-order,
interatomic force constants (IFC2), the third-order coefficients
define the cubic, or third-order, force constants (IFC3).
The harmonic approximation retains
\cref{eq:e_exp_ph} to second order, and is sufficient for phonon
dispersion and most equilibrium thermodynamic quantities
\cite{dove_introduction_1993, wallace_thermodynamics_1998}.
Anharmonic lattice-dynamics methods include the cubic and potentially higher-order terms.
These introduce phonon--phonon scattering processes and hence the finite linewidth,
frequency shift, and thermal transport properties of crystalline solids
\cite{ziman_electrons_2001, srivastava_physics_2019}.

% ----------------------------------------------------------------------
\subsection{Harmonic force constants (IFC2)}
\label{sub:fc2}
% ----------------------------------------------------------------------

Introducing the primitive-cell degree-of-freedom (DOF) index
$\mu = (\kappa, \alpha)$, which bundles the basis atom $\kappa$ and
the Cartesian direction $\alpha$ into a single label, the harmonic
force constants are defined as
\begin{equation}\label{eq:fc2_def}
  \Phi(0\mu,\, l'\mu')
  = \frac{\partial^2 E}
         {\partial u_{0\mu}\,\partial u_{l'\mu'}},
\end{equation}
where $u_{l\mu}$ is the displacement of DOF $\mu$ in cell $l$.
Within a perfect crystal, translational symmetry lets us fix the
first cell index to the home cell ($l{=}0$) without loss of
generality, so that $\Phi(0\mu,\,l'\mu')$ depends only on the
relative cell separation.  Grouping the three Cartesian directions
for a fixed pair of atoms $(\kappa,\kappa')$ yields a
$3\times 3$ block whose rows and columns correspond to force and
displacement components, respectively.

The IFC2 quantify how a small displacement of DOF $\mu'$ in cell
$l'$ changes the force on DOF $\mu$ in the home cell.  Linearising
the Newtonian equations of motion around equilibrium then gives the
force--displacement relation
\begin{equation}\label{eq:force-disp}
  F(0\mu)
  = -\sum_{l'\mu'} \Phi(0\mu,\, l'\mu')\, u_{l'\mu'},
\end{equation}
which is crucial to both harmonic lattice dynamics and the
finite-displacement schemes used to compute the IFC2 discussed in
\cref{sub:fc_computation}.

By construction, the IFC2 tensor obeys three exact relations that are
useful both for reducing storage and for enforcing physical
consistency when force constants are obtained numerically.
The first is permutation symmetry: since partial derivatives
commute, the tensor is symmetric under exchange of its two legs,
\begin{equation}\label{eq:fc2_perm}
  \Phi(0\mu,\,l'\mu') = \Phi(l'\mu',\,0\mu),
\end{equation}
which, combined with translational invariance, can equivalently be
written as
\begin{equation}
  \Phi(0\mu,\,l'\mu')
  =
  \Phi(0\mu',\,(-l')\mu).
  \label{eq:ifc-symmetry-equivalent}
\end{equation}
This form is commonly used in lattice-dynamical treatments of force
constants
\cite{born_dynamical_1996, maradudin_theory_1963, gonze_dynamical_1997}.

The second symmetry is translational invariance: a rigid translation of
the whole crystal in any Cartesian direction $\alpha'$ cannot generate
a restoring force, which leads to the acoustic sum rule
\begin{equation}\label{eq:fc2_asr}
  \sum_{l'\kappa'} \Phi_{\alpha\alpha'}(0\kappa,\,l'\kappa') = 0
  \qquad \forall\;\kappa,\alpha,\alpha',
\end{equation}
where the condition must hold for each triple
$(\kappa,\alpha,\alpha')$ separately \cite{born_dynamical_1996, pick_microscopic_1970, gonze_dynamical_1997}.
This relation guarantees that the three
acoustic branches have zero frequency at the $\Gamma$
point, and in practice its enforcement is critical for numerically
computed IFC2, as even small residual violations produce spurious
acoustic gaps in the dispersion.  The third is space-group
symmetry: under each operation $\{S|\vc{v}\}$ of the crystal point
group, the IFC2 transform as a second-rank Cartesian tensor on the
index pair and simultaneously map one atom--cell pair onto another.
The resulting linear relations among the entries of $\Phi$
typically reduce the number of independent parameters by one or two
orders of magnitude, and are exploited systematically by software
packages such as phonopy \cite{togo_first_2015} and symfc
\cite{seko_projector-based_2024}, and are discussed in-depth in \citet{fu_group_2019}.

Phonon frequencies and polarisation vectors are obtained as
eigenpairs of the mass-weighted Fourier transform of $\Phi$, the
dynamical matrix (see \cref{sec:notation_fourier} for the two phase
conventions used in this work).  See the standard literature, e.g.
\citet{born_dynamical_1996}, \citet{maradudin_theory_1963}, and
\citet{dove_introduction_1993}.

% ----------------------------------------------------------------------
\subsection{Anharmonic force constants (IFC3)}
\label{sub:fc3}
% ----------------------------------------------------------------------

The cubic, or third-order, interatomic force constants are the
generalisation of \cref{eq:fc2_def} to the next order of the
Taylor expansion \cref{eq:e_exp_ph}.  In the composite-index
notation, they read
\begin{equation}\label{eq:fc3_def}
  \Phi^{(3)}(0\mu,\,l'\mu',\,l''\mu'')
  = \frac{\partial^3 E}
         {\partial u_{0\mu}\,
          \partial u_{l'\mu'}\,
          \partial u_{l''\mu''}},
\end{equation}
and, as before, translational symmetry lets us fix the first cell
index to the home cell.  The IFC3 tensor captures the leading
anharmonic corrections to lattice dynamics.  It sets the amplitude
of three-phonon scattering processes, determines the finite
linewidth and temperature-dependent frequency shift of phonon modes,
and provides the dominant resistive mechanism for lattice thermal
conductivity in most crystalline solids \cite{ziman_electrons_2001, srivastava_physics_2019}.

The symmetry properties of IFC3 parallel those of IFC2, but with a
richer permutation structure due to the three legs of the
tensor.  Commutativity of partial derivatives now gives full $S_3$
symmetry, so that $\Phi^{(3)}(0\mu,\,l'\mu',\,l''\mu'')$ is
invariant under all six simultaneous permutations of its three
cell--DOF pairs:
\begin{equation}\label{eq:fc3_perm}
  \Phi^{(3)}(0\mu,\,l'\mu',\,l''\mu'')
  = \Phi^{(3)}(l'\mu',\,0\mu,\,l''\mu'')
  = \Phi^{(3)}(l''\mu'',\,l'\mu',\,0\mu)
  = \cdots
\end{equation}
\cite{srivastava_physics_2019}.  Translational
invariance again generates acoustic sum rules, now one per leg:
summing $\Phi^{(3)}$ over the cell--atom labels of any one leg must
vanish for all choices of the remaining indices,
\begin{equation}\label{eq:fc3_asr}
  \sum_{l''\kappa''}
     \Phi^{(3)}_{\alpha\alpha'\alpha''}(0\kappa,\,l'\kappa',\,l''\kappa'')
  = 0
  \qquad \forall\;\kappa,\kappa',l',\alpha,\alpha',\alpha'',
\end{equation}
and analogously under summation over the first or second leg
\cite{srivastava_physics_2019}.  As for the
harmonic case, enforcing these sum rules on numerically extracted
IFC3 is essential for a consistent phonon self-energy and, hence,
for accurate thermal-transport predictions \cite{li_shengbte_2014}.
Space-group symmetry acts as a third-rank Cartesian tensor on the
index triple while simultaneously permuting the atom--cell labels,
and typically reduces the independent-entry count by one to two
orders of magnitude \cite{fu_group_2019, seko_projector-based_2024,
esfarjani_method_2008}.

Despite these reductions, the storage requirements for the full IFC3
grow quickly with system size.  After fixing one cell index by
translational symmetry, the number of independent entries scales as
$\cO(\nprim\,\ndof^2) = \cO(\Ncell^2\nprim^3)$, which rapidly
becomes prohibitive for large supercells.  This scaling,
together with the even more unfavourable
cost of the four-tensor contraction that enters the phonon--phonon
self-energy, motivates the compression schemes developed in
\cref{sub:fc3_compression}.

% ----------------------------------------------------------------------
\subsection{Computation of force constants}
\label{sub:fc_computation}
% ----------------------------------------------------------------------

Several established methods exist for obtaining interatomic force
constants from first-principles calculations.  We summarise the
main approaches below. The choice among them involves trade-offs
between computational cost, accessible interaction range, the
order of anharmonicity that can be treated, and availability of software.

% ----------------------------------------------------------------------
\subsubsection{Finite displacements}
\label{ssub:fd_fc}
% ----------------------------------------------------------------------
The most widely used approach constructs a supercell, displaces atoms
one at a time along symmetry-inequivalent directions, performs a
self-consistent DFT calculation for each displaced structure, and
reconstructs the force-constant tensor by inverting the
force--displacement relation \cref{eq:force-disp}.  The method was
introduced by \citet{parlinski_first-principles_1997} and is
implemented in codes such as phonopy \cite{togo_first_2015},
PHON \cite{alfe_phon_2009}, and Phonon
\cite{parlinski_first-principles_1997}.  Space-group symmetry reduces
the number of independent displacements and hence the number of
required DFT calculations, making the approach practical for cells of
moderate size.

For cubic force constants, the analogous procedure requires displacing
two atoms simultaneously. The number of required supercell
calculations therefore grows quadratically with the number of atoms in
the primitive cell
\cite{esfarjani_method_2008, li_shengbte_2014}.  \Citet{esfarjani_method_2008}
introduced a general method that extracts harmonic and anharmonic FCs
to arbitrary neighbour shell from a single set of displacements by
solving a linear least-squares problem with symmetry constraints.
Widely used implementations for the IFC3 include thirdorder.py
(distributed with ShengBTE) \cite{li_shengbte_2014} and phono3py
\cite{togo_implementation_2023}.  The
latter also provides the infrastructure for computing phonon
linewidths, thermal conductivity, and spectral functions from the
extracted IFC3.

\paragraph{Supercell convergence and aliasing.}
The choice of supercell determines the real-space range over which the
force constants can be resolved before periodic wrap-around introduces
aliasing.  For a supercell of linear dimension $N$ along a lattice
direction, the calculated force constant at separation $r$ is the
periodically folded sum
\begin{equation}\label{eq:fc_aliasing}
  \Phi_{\mathrm{per}}^{(N)}(r)
  = \sum_{p\in\mathbb{Z}} \Phi(r + pN).
\end{equation}
In a $2\times 2\times 2$ supercell, the resolvable relative
translations are limited to those representable within the finite
periodic cell -- contributions from more distant cells are folded back
onto these separations.  A larger supercell therefore allows force
constants to be resolved over a wider real-space range before aliasing
occurs \cite[Chs.~2,~6]{dove_introduction_1993}.

This aliasing directly affects the phonon dispersion. At wavevectors
$\vq$ that are not commensurate with the supercell mesh, the
Fourier-interpolated dynamical matrix inherits errors from the
folded-back tails of $\Phi$. In practice, convergence must therefore be
checked for the target observable---phonon frequencies, thermal
conductivity, or self-energy---by increasing the supercell size until
the aliased contributions are negligible
\cite{togo_first_2015, parlinski_first-principles_1997}.

For the IFC3, the aliasing issue is compounded by the fact that three
cell indices are involved: each pair in the triplet
$(0\kappa, l'\kappa', l''\kappa'')$ must lie within the
interaction range resolvable by the supercell.  This makes third-order
properties more demanding in supercell size than their harmonic
counterparts \cite{li_shengbte_2014, esfarjani_method_2008}.

% ----------------------------------------------------------------------
\subsubsection{Density-functional perturbation theory (DFPT)}
\label{ssub:dfpt_fc}
% ----------------------------------------------------------------------
Density functional perturbation theory avoids explicit supercell
calculations by computing the linear response of the electron density
to a phonon perturbation at an arbitrary wave-vector $\vq$ within the
primitive cell.  The method was introduced by
\citet{baroni_greens-function_1987} and
\citet{gonze_adiabatic_1995}, with a comprehensive review by
\citet{baroni_phonons_2001}.  Second-order force constants at any
$\vq$ are obtained directly from the self-consistent solution of the
Sternheimer equations, and real-space IFC2 are obtained by Fourier
interpolation of the dynamical matrix computed on a regular mesh of
$\vq$-points \cite{gonze_interatomic_1994}.

For cubic force constants, the $2n{+}1$ theorem allows $\Phi^{(3)}$
to be extracted from the first-order wavefunction and density
responses already computed at the IFC2 stage, without requiring
second-order responses \cite{debernardi_anharmonic_1995, debernardi_third-order_1994}.
\citet{deinzer_ab_2003} applied this
to compute phonon linewidths in Ge and Si, and
\citet{paulatto_anharmonic_2013} generalised it to arbitrary
triplets $(\vq_1,\vq_2,\vq_3)$ with $\vq_1{+}\vq_2{+}\vq_3 =\vc{G}$,
enabling the direct evaluation of three-phonon scattering matrix
elements throughout the Brillouin zone without real-space
supercells.  \citet{paulatto_first-principles_2015} later extended
the approach to compute phonon spectral functions with strong anharmonicity.

DFPT-based codes differ substantially in what they support at the
IFC2 and IFC3 levels.  For harmonic IFCs, Quantum ESPRESSO
\cite{giannozzi_quantum_2009, giannozzi_advanced_2017} (via
\texttt{ph.x}) supports norm-conserving (NC), ultrasoft (USPP), and
projector-augmented-wave (PAW) pseudopotentials, with LDA, GGA, and
(in recent versions) hybrid functionals. ABINIT
\cite{gonze_abinit_2009} similarly supports both NC and PAW datasets,
non-magnetic, collinear, and non-collinear spin, as well as DFT$+U$.
VASP \cite{kresse_efficient_1996} provides a more rudimentary
DFPT implementation (\texttt{IBRION}\,{=}\,7,\,8) that is restricted
to displacements commensurate with the supercell ($\vq{=}\vc 0$) and
to LDA/GGA functionals.

For cubic force constants, the picture is more restrictive.  In
Quantum ESPRESSO, $\Phi^{(3)}$ is available through the D3Q and
\texttt{thermal2} packages
\cite{paulatto_anharmonic_2013, fugallo_ab_2013}.  These
codes support only norm-conserving pseudopotentials and only LDA or
GGA exchange--correlation functionals. Neither USPP, PAW, meta-GGA,
hybrid functionals, nor DFT$+U$ are currently implemented at third
order.  The same limitation applies to all third-order energy
derivatives in Quantum ESPRESSO, including Raman cross-sections.  In
ABINIT, a framework for third-order derivatives of the total energy
(3DTE) exists and is used in literature for non-linear optical
responses and Raman tensors \cite{veithen_nonlinear_2005, romero_abinit_2020}.
The extraction of the full three-phonon IFC3
tensor on a $\vq$-grid analogous to D3Q, however, is not available and, in
practice, anharmonic lattice dynamics in ABINIT is usually performed via
the supercell-based a-TDEP module \cite{bottin_-tdep_2020} rather than through pure
DFPT.  VASP does not currently implement third-order DFPT, and cubic
force constants with VASP are typically obtained through
finite--displacement calculations with, e.g., phono3py
\cite{togo_distributions_2015}. Hence, when computing IFC3 with DFPT, one is
effectively restricted to a LDA/GGA exchange correlation with
norm-conserving pseudopotentials, via Quantum ESPRESSO + D3Q.

% ----------------------------------------------------------------------
\subsubsection{Compressive sensing and machine-learning approaches}
\label{ssub:csld_fc}
% ----------------------------------------------------------------------
Compressive-sensing lattice dynamics (CSLD) casts the force-constant
extraction as an $\ell_1$-regularised regression on DFT forces
computed for a small number of randomly displaced configurations
\cite{zhou_lattice_2014, zhou_compressive_2019}.  The sparsity prior
selects the minimal set of non-negligible FC entries, potentially
capturing long-range couplings that would require prohibitively large
supercells in the standard finite-displacement scheme.  The hiPhive
package \cite{eriksson_hiphive_2019} implements this idea and uses
crystal-symmetry reduction and advanced regression
algorithms to extract force constants up to arbitrary order
from semi-random displacement configurations. These regression-based approaches are
particularly attractive when the full enumeration of finite
displacements becomes prohibitively expensive for large systems.

% ----------------------------------------------------------------------
\subsection{FC3 Compression}
\label{sub:fc3_compression}
% ----------------------------------------------------------------------

Storing the full FC3 tensor requires $\cO(\nprim\,\ndof^{2}) = \cO(\Ncell^{2}\nprim^{3})$
memory, which is prohibitive for large-scale simulations. Beyond storage, the dense four-tensor
contraction entering the phonon--phonon self-energy scales as $\cO(\ndof^{4})$ per frequency
point and dominates the cost of the quantum-transport simulation. It is hence desirable to
efficiently approximate the FC3 interactions.

Effectively every published anharmonic phonon NEGF implementation avoids the problem by
dropping self-energy and FC3 entries at large atomic separations
\cite{guo_quantum_2020, guo_anharmonic_2021, luisier_atomistic_2012, li_anharmonic_2009,
      polanco_nonequilibrium_2021, mingo_anharmonic_2006}.
The standard choice is a real-space pair-distance cut-off: keep only those triplets
$(0\kappa, l'\kappa', l''\kappa'')$ whose largest pairwise distance lies below some
$r_\mathrm{c}$. The same idea is used in the BTE workflows of ShengBTE and its successors,
which restrict FC3 to a few nearest-neighbour shells \cite{li_shengbte_2014}.

The drawback is that the approximation error is uncontrolled: there is no variational bound
on what is discarded, and the truncation removes genuine non-local contributions that can
matter quantitatively in the phonon self-energy. In anharmonic NEGF this is already visible in
silicon. \citeauthor{guo_quantum_2020} use first-principles FC3 in a quantum-transport treatment
of cross-plane heat flow in Si thin films, where they compare first- and third-nearest-neighbour
FC3 cutoffs \cite{guo_quantum_2020}. Using the same FC3 data in ShengBTE, they report bulk Si
thermal conductivities at 300~K of $120.69$, $136.46$, and $147.13~\mathrm{W\,m^{-1}K^{-1}}$ when
FC3 was truncated at the first, second, and third neighbour shell, respectively, so that a
first-shell cut-off underestimates the converged bulk value by nearly $20\%$. For bulk silicon,
including up to 4th nearest neighbour interactions gives converged results
\cite{li_shengbte_2014, qin_accelerating_2018}.

For three-dimensional covalent crystals, the cut-off problem is usually less severe than in
two-dimensional materials. \citeauthor{mcgaughey_phonon_2019} report heat conductivities
$\kappa_\ell(300\,\mathrm{K}) = 150, 152, 155, 151,$ and $149~\mathrm{W\,m^{-1}K^{-1}}$ for
cubic cut-offs of $2$, $3$, $4$, $5$, and $6$ neighbour shells, respectively in bulk Si, indicating
that once a few shells are retained the residual cut-off dependence is already small
\cite{mcgaughey_phonon_2019}. \citeauthor{jiang_how_2022} reached the same qualitative conclusion
by contrasting graphene and silicene against bulk Si: whereas the 2D systems require inclusion of
up to 14th order NN shells, bulk Si shows ``no convergence issue'' even for very small neighbour
cut-offs \cite{jiang_how_2022}.

That behaviour is, however, not universal across 3D materials. \citeauthor{qin_accelerating_2018}
showed that bulk SnSe still exhibits strong interactions at distances around $6.15~\text{\AA}$,
and that these coincide with a sharp change in the converged $\kappa_\ell$ once the pair cut-off
exceeds roughly $6~\text{\AA}$ \cite{qin_accelerating_2018}. They also show that for the 2D
materials studied, the observables in BTE diverge for increasing $q$-grid when the cut-off is too
small.

We therefore seek a representation of the FC3 tensor that
\begin{itemize}
    \item provides a controllable approximation error,
    \item achieves significant compression of the stored data, and
    \item turns the dense $\cO(\ndof^{4})$ self-energy contraction into a sum of low-order
          contractions whose cost scales polynomially in a parameter~$R$.
\end{itemize}
A large body of work has developed such representations in similar settings.
In quantum chemistry, tensor hypercontraction (THC) \cite{hohenstein_tensor_2012,
parrish_tensor_2012} and density fitting \cite{whitten_coulombic_1973} factorise the four-index
electron-repulsion integral tensor into products of matrices. In many-body physics, tensor
networks (matrix product states, tensor trains) \cite{oseledets_tensor-train_2011,
cichocki_tensor_2016} compress high-dimensional objects at controlled error. For the present
setting, lattice dynamics, tensor decomposition has recently proved effective:
\citeauthor{luo_data-driven_2024} used a truncated SVD on the electron-phonon matrix in Wannier
basis and recovered high-accuracy transport, spin, and superconducting observables from $1$--$2\%$
of the singular values \cite{luo_data-driven_2024}. The same group then proposed the permanent
CP (PCP) ansatz for the $n$-IFC tensors and reported $10^{3}$--$10^{4}$ compression factors
across Si, HgTe, MgO, TiNiSn and monoclinic ZrO$_2$ while retaining $>98\%$ accuracy on
$\kappa_\ell$ for three- and four-phonon scattering \cite{luo_tensor_2025}. Compressive-sensing
lattice dynamics (CSLD) is a method applied to a related problem at the fitting stage: an
$\ell_1$-regularised regression on DFT forces selects a sparse set of non-negligible FC entries
directly \cite{zhou_lattice_2014, zhou_compressive_2019}. Projector-based approaches, implemented
in the symfc software, construct a complete orthonormal basis for FCs compatible with all
crystal symmetries \cite{seko_projector-based_2024}, while group-theoretical reductions enumerate
the independent FC3 entries from the space group \cite{fu_group_2019}. In the following, we take
the fully-populated FC3 tensor as given (from finite displacements, DFPT, or CSLD) and compress
its representation. However, one can likely also fit a compressed tensor from DFT data with an
adaption of CSLD or similar methods without ever obtaining a full FC3 tensor.

\citet{kolda_tensor_2009} organise tensor compression into two principal families: the
Tucker decomposition, which writes a tensor as a core tensor multiplied by a factor matrix
along each mode (generalising principal-component analysis), and the CP
(CANDECOMP/PARAFAC) decomposition, which writes it as a sum of rank-one outer products
(generalising the matrix SVD in a different direction). Applied to a tensor with the FC3 index
pattern and combined with the possible symmetry constraints on the factor matrices, these two
families yield five canonical ans\"atze, ordered here from least to most
tensor-structure-aware:
\begin{enumerate}
  \item Truncated matricization SVD, i.e.\ the classical matrix SVD applied to an matricization of
        the tensor. fully algebraic and Eckart--Young optimal in the chosen matricization, but
        separating only one leg from the other two.
  \item Truncated HOSVD (Tucker), with independent factor matrices per mode and a small dense core.
  \item CP, obtained by restricting the Tucker core to be diagonal, with three
        independent rank-$R$ factor matrices and no symmetry constraint.
  \item INDSCAL \cite{carroll_analysis_1970}, obtained from CP by setting the
        factor matrices on the two internal legs equal, exploiting the
        $S_2$ symmetry $\Phi_{\mu j k} = \Phi_{\mu k j}$.
  \item Symmetric CP, equivalent to the Waring decomposition of a homogeneous
        cubic polynomial \cite{brachat_symmetric_2010, comon_symmetric_2008,
        oeding_eigenvectors_2013}, obtained from CP by setting all three factor
        matrices equal (using the $p(\mu)$ lift of a primitive-cell DOF to its
        supercell counterpart in cell zero), exploiting the full $S_3$
        symmetry.
\end{enumerate}
Entries 1 and 2 are non-iterative. Entries 3--5 are NP-hard in the worst case
\cite{hillar_most_2013} and use iterative optimisation, with the known caveats of
rank-degeneracy and non-closedness of rank-$R$ sets
\cite{stegeman_degeneracy_2006, de_silva_tensor_2008}. Each added symmetry constraint from
entry~2 onwards reduces the parameter count at matched rank while shrinking the set of
representable tensors. As the transport-residual study of \cref{sec:res_fc}
shows, the most-constrained ansatz is not necessarily the least accurate per
parameter --- the transport observable converges faster than the Frobenius
residual the fitters minimise.

% ----------------------------------------------------------------------
\subsubsection{Truncated matricization SVD}
\label{sub:fc3_matsvd}
% ----------------------------------------------------------------------
Stack the $3\nprim$ slices $M_\mu = \Phi_{\mu,:,:}\in\mathbb{R}^{\ndof\times\ndof}$
into a single matrix
\begin{equation}\label{eq:fc3_matsvd_stack}
  M_{\mathrm{stack}}
  = \begin{pmatrix} M_1 \\ \vdots \\ M_{3\nprim} \end{pmatrix}
  \in \mathbb{R}^{3\nprim\,\ndof \,\times\, \ndof}.
\end{equation}
A truncated matrix SVD at rank $R$ gives
\begin{equation}\label{eq:fc3_matsvd_ansatz}
    \tilde\Phi_{\mu j k}^{\mathrm{mSVD}}
    = \sum_{r=1}^{R} F_r(\mu,j)\; h_r(k),
\end{equation}
with $F_r(\mu,j) = \sigma_r\,[\vc{u}_r]_{(\mu-1)\ndof+j}$ and $h_r(k)=[\vc{v}_r]_k$, where
$\sigma_r$, $\vc{u}_r$, $\vc{v}_r$ are the $R$ leading singular values, left singular vectors,
and right singular vectors of $M_{\mathrm{stack}}$. This is equivalent to a Tucker1 decomposition
with compression on mode~3 only \cite[\S4]{kolda_tensor_2009}.

\paragraph{Computation.} A single matrix SVD of $M_{\mathrm{stack}}$, computed by
standard dense LAPACK routines or, for large $\ndof$, by randomised linear algebra algorithms
\cite{halko_finding_2011}. The Eckart--Young--Mirsky theorem
\cite{eckart_approximation_1936, mirsky_symmetric_1960} guarantees that, at the chosen $(\mu,j)|k$ bipartition,
the truncation is globally optimal among all rank-$R$ approximations in the Frobenius norm.
This makes it the natural baseline against which the structurally richer decompositions of
\cref{sub:fc3_tucker,sub:fc3_cp,sub:fc3_indscal,sub:fc3_symcp} must be compared. The ansatz does
not exploit the $S_2$ symmetry on $(j,k)$: modes~1 and~2 are folded into a single flat index, so
parameters are wasted on storing redundant information and the reconstruction is not exactly
symmetric in $(j,k)$ unless restored in post-processing.

\paragraph{Parameter count.}
\begin{equation}\label{eq:fc3_matsvd_params}
  N_{\mathrm{param}}^{\mathrm{mSVD}}(R)
  = R\bigl((3\nprim+1)\,\ndof + 1\bigr)
  \;\sim\; \cO\!\bigl(R\,\nprim\,\ndof\bigr)
  \;=\; \cO\!\bigl(R\,\Ncell\,\nprim^{2}\bigr),
\end{equation}
one factor of $\nprim$ worse than the CP-family ans\"atze below.

\paragraph{Self-energy kernel.} The separable structure factorises the ring contraction entering
the phonon--phonon self-energy into two independent bilinear forms. The $\vq$-convolution is an
FFT over the $R^{2}$ rank pairs, replacing the dense $\ndof^{4}$ contraction by
$R^{2}\ndof$-scaling accumulations.

% ----------------------------------------------------------------------
\subsubsection{Truncated HOSVD (Tucker decomposition)}
\label{sub:fc3_tucker}
% ----------------------------------------------------------------------
The Tucker ansatz writes FC3 as a multilinear contraction of a small core tensor with one factor
matrix per mode \cite{tucker_mathematical_1966, kolda_tensor_2009}:
\begin{equation}\label{eq:fc3_tucker_ansatz}
    \tilde\Phi_{\mu j k}^{\mathrm{HOSVD}}
    = \sum_{p=1}^{R_1}\sum_{q=1}^{R_2}\sum_{r=1}^{R_2}
      G_{p q r}\; A_{\mu p}\; B_{j q}\; B_{k r},
\end{equation}
with $A\in\mathbb{R}^{3\nprim\times R_1}$, $B\in\mathbb{R}^{\ndof\times R_2}$ columnwise
orthonormal, and a core $G\in\mathbb{R}^{R_1\times R_2\times R_2}$. To preserve the $S_2$
symmetry of the two internal legs, the factor matrices on modes~2 and~3 can be chosen identical ($B=C$)
and the core is constrained to satisfy $G_{pqr}=G_{prq}$.

\paragraph{Computation.} The truncated higher-order SVD (HOSVD) of
\citet{de_lathauwer_multilinear_2000} produces $A$ and $B$ in closed form as the $R_1$ and $R_2$
leading left singular vectors of the mode-$1$ and mode-$2$ unfoldings of $\Phi$, followed by the
core contraction $G = \Phi\times_1 A^\top\times_2 B^\top\times_3 B^\top$ (which is automatically
$S_2$-symmetric when $\Phi$ is). This is quasi-optimal: the truncation error is bounded by a
factor $\sqrt{3}$ above the best rank-$(R_1,R_2,R_2)$ approximation error
\cite{hackbusch_tensor_2019, oseledets_tensor-train_2011}. The fit can optionally be refined by a
few steps of higher-order orthogonal iteration (HOOI), an alternating least squares (ALS)-type loop
 \cite{de_lathauwer_best_2000}. The special case $R_1=3\nprim$, $A=I$ (compression of the two
internal legs only) is the Tucker1 baseline and is obtained by a single SVD of the corresponding
unfolding.

\paragraph{Parameter count.}
\begin{equation}\label{eq:fc3_tucker_params}
  N_{\mathrm{param}}^{\mathrm{HOSVD}}(R_1,R_2)
  = 3\nprim\,R_1 + \ndof\,R_2 + R_1\,\tfrac{R_2(R_2+1)}{2}
  \;\sim\; \cO\!\bigl(R_2\,\ndof + R_1\,R_2^{2}\bigr),
\end{equation}
where the $S_2$ constraint halves the nominal core cost $R_1 R_2^{2}$ to $R_1 R_2(R_2+1)/2$.
The $R_2^{2}$-scaling core term makes Tucker the most parameter-heavy of the five ans\"atze
when $R_2$ is not small compared to $\ndof^{1/2}$. for strongly anharmonic systems where the
effective rank is moderate this is mitigated by allowing $R_1\ll R_2$.

\paragraph{Self-energy kernel.} The self-energy contraction factorises into three mode-wise
contractions against $A$ and $B$, leaving the $R_1\times R_2\times R_2$ core as a small dense
sub-problem. The total cost scales as $R_2^{2}(R_1+\ndof)$ per frequency point rather than
$\ndof^{4}$.

% ----------------------------------------------------------------------
\subsubsection{Canonical Polyadic decomposition (CP)}
\label{sub:fc3_cp}
% ----------------------------------------------------------------------
Dropping the off-diagonal entries of the Tucker core yields the CP approximation, also
known as CANDECOMP and PARAFAC \cite{harshman_foundations_1970, hitchcock_expression_1927, carroll_analysis_1970}, a sum of
rank-one outer products with three independent factor vectors per term:
\begin{equation}\label{eq:fc3_cp_ansatz}
    \tilde\Phi_{\mu j k}^{\mathrm{CP}}
    = \sum_{r=1}^{R} \lambda_r\;a_r(\mu)\;b_r(j)\;c_r(k).
\end{equation}
No symmetry between the modes is imposed. For a tensor that is
$S_2$- or $S_3$-symmetric, the fit does not necessarily converge to a symmetric
solution \cite{ten_berge_typical_2004}. Uniqueness can be shown under certain conditions
\cite{kruskal_three-way_1977, sidiropoulos_uniqueness_2000} and is generically
guaranteed for the FC3 tensor of interest here
\cite{chiantini_generic_2012}.

\paragraph{Computation.} The standard algorithm is alternating least squares
(CP-ALS), which cycles through the three modes solving, at each step, a
linear least-squares problem for one factor matrix while the other two are
fixed \cite[Sec.~3.4]{kolda_tensor_2009}. Convergence is only to a
stationary point and can be slow. Variants include
damped Gauss--Newton (dGN), PMF3 \cite{paatero_weighted_1997},
line-search extrapolation (ELS) \cite{rajih_enhanced_2008}, and
simultaneous-diagonalisation approaches \cite{de_lathauwer_link_2006}
that recast CP as a generalised eigenproblem when
$R\le\min(3\nprim,\ndof)$.

\paragraph{Parameter count.}
\begin{equation}\label{eq:fc3_cp_params}
  N_{\mathrm{param}}^{\mathrm{CP}}(R)
  = R\,\bigl(3\nprim + 2\,\ndof + 1\bigr)
  \;\sim\; \cO\!\bigl(R\,\ndof\bigr)
  \;=\; \cO\!\bigl(R\,\Ncell\,\nprim\bigr).
\end{equation}

% ----------------------------------------------------------------------
\subsubsection{INDSCAL (CP with internal-leg symmetry)}
\label{sub:fc3_indscal}
% ----------------------------------------------------------------------
Individual-differences-in-scaling (INDSCAL) \cite{carroll_analysis_1970, kolda_tensor_2009} is the
specialisation of CP in which the factor matrices on two modes are constrained equal, tailored
to tensors whose slices $\Phi_{\mu,:,:}$ are symmetric:
\begin{equation}\label{eq:fc3_indscal_ansatz}
    \tilde\Phi_{\mu j k}^{\mathrm{INDSCAL}}
    = \sum_{r=1}^{R} d_r(\mu)\; v_r(j)\; v_r(k),
\end{equation}
with $d_r\in\mathbb{R}^{3\nprim}$ and $v_r\in\mathbb{R}^{\ndof}$. This imposes the $S_2$
symmetry $\Phi_{\mu j k}=\Phi_{\mu k j}$. Typical and maximal ranks for
partially symmetric tensors were characterised by
\citet{ten_berge_typical_2004}: for the shapes relevant here, they coincide with or
are strictly smaller than the ranks of generic nonsymmetric tensors of the same shape, so the
symmetry constraint is not costly in expressive power.

\paragraph{Computation.} \citet{carroll_analysis_1970} originally proposed running unconstrained
CP-ALS with the two internal factors $B$ and $C$ updated separately, relying on symmetry of the
data to drive $B\to C$. \citet{ten_berge_typical_2004, ten_berge_carroll_2009} showed that this convergence is generic
but not guaranteed. Modern implementations impose equality explicitly via a symmetric ALS or via
a Gauss--Newton method \cite{kolda_tensor_2009, kolda_numerical_2015}.
A cheap closed-form alternative performs a matrix SVD of $\Phi$ reshaped to
$3\nprim\times\ndof^{2}$ followed by an eigendecomposition of each symmetric slice of the
resulting right factors, and keeps the $R$ terms with largest amplitude. This is not jointly
optimal (the total rank needed to match a given error is typically larger than that of a
jointly-fitted INDSCAL) but is useful as an initialisation.

\paragraph{Parameter count.}
\begin{equation}\label{eq:fc3_indscal_params}
  N_{\mathrm{param}}^{\mathrm{INDSCAL}}(R)
  = R\,\bigl(3\nprim + \ndof\bigr)
  \;\sim\; \cO\!\bigl(R\,\ndof\bigr)
  \;=\; \cO\!\bigl(R\,\Ncell\,\nprim\bigr),
\end{equation}
smaller than CP by one $\ndof$-length factor per term (a constant $\sim2$, not a
factor of $\ndof$), from merging $b_r\equiv c_r$.

\paragraph{Self-energy kernel.} Identical in scaling to CP, $\cO(R^{2}\ndof)$ per frequency, but
with the additional property that the reconstruction is exactly $S_2$-symmetric in $(j,k)$, so
ASR-sensitive observables receive only the quadratic-in-error contribution
(\cref{sec:res_fc}).


% ----------------------------------------------------------------------
\subsubsection{Symmetric CP (Waring) decomposition}
\label{sub:fc3_symcp}
% ----------------------------------------------------------------------
The symmetric CP decomposition imposes equality of all three factor matrices, $A=B=C$, on the
fully $S_3$-symmetrised FC3 tensor obtained via the supercell lift
$p:\{1,\dots,3\nprim\}\to\{1,\dots,\ndof\}$. The Waring parametrisation reads:
\begin{equation}\label{eq:fc3_waring_ansatz}
    \tilde\Phi_{\mu j k}^{\mathrm{Waring}}
    = \sum_{r=1}^{R} \lambda_r\;
      v_r\!\bigl(p(\mu)\bigr)\; v_r(j)\; v_r(k),
\end{equation}
with $v_r\in\mathbb{R}^{\ndof}$ and $\lambda_r\in\mathbb{R}$. This is the tensor form of the
classical Waring problem: expressing a homogeneous cubic polynomial as a sum of $R$ cubes of
linear forms \cite{brachat_symmetric_2010, comon_symmetric_2008, oeding_eigenvectors_2013}. For
generic third-order symmetric tensors on $\ndof$ indices, the Waring rank equals
$\lceil\binom{\ndof+2}{3}/\ndof\rceil$ by the Alexander--Hirschowitz theorem
\cite{alexander_generic_1997} (which is a non-trivial statement and proven by Alexander and Hirschowitz
in a set of 5 papers of over 100 pages). An alternative ``permanent-CP'' parametrisation, introduced
by \citet{luo_tensor_2025} as PCP, uses three distinct vectors per term explicitly symmetrised
over $S_3$,
\begin{equation}\label{eq:fc3_pcp_ansatz}
    \tilde\Phi_{\mu j k}^{\mathrm{PCP}}
    = \sum_{\xi=1}^{R} \frac{\lambda_\xi}{6}
      \sum_{\sigma\in S_3}
      u^{\xi}_{\sigma_1}\!\bigl(p(\mu)\bigr)\;
      u^{\xi}_{\sigma_2}(j)\;
      u^{\xi}_{\sigma_3}(k),
\end{equation}
with $u^{\xi}_1, u^{\xi}_2, u^{\xi}_3\in\mathbb{R}^{\ndof}$. PCP fits the same class of symmetric
tensors as Waring but with more capacity per term (at threefold parameter cost) and a smoother
optimisation landscape.

\paragraph{Computation.} Algorithms fall into two categories. Algebraic methods reconstruct the
decomposition from moment and flattening matrices via generalised eigenproblems and are exact up
to numerical precision when the rank is below the generic bound
\cite{brachat_symmetric_2010, oeding_eigenvectors_2013}. Numerical methods optimise
$\|\Phi-\tilde\Phi\|_F^{2}$ directly: the shifted symmetric higher-order power method (SS-HOPM)
\cite{kolda_shifted_2011} and the subspace power method \cite{kileel_subspace_2025} extract
components one at a time with provable convergence for low-rank tensors; all-at-once gradient
and Gauss--Newton methods \cite{kolda_numerical_2015} fit all $R$ components simultaneously.
Degeneracy and diverging-norm instabilities can appear because the set of symmetric tensors of
rank $\le R$ is not closed in general
\cite{stegeman_degeneracy_2006, de_silva_tensor_2008}. Remedies include norm penalties and
second-order line searches.

\paragraph{Parameter count.}
\begin{align}
  N_{\mathrm{param}}^{\mathrm{Waring}}(R) &= R\,(\ndof + 1)
    \;\sim\; \cO\!\bigl(R\,\ndof\bigr)
    \;=\; \cO\!\bigl(R\,\Ncell\,\nprim\bigr),
    \label{eq:fc3_waring_params}\\
  N_{\mathrm{param}}^{\mathrm{PCP}}(R) &= R\,(3\,\ndof + 1)
    \;\sim\; \cO\!\bigl(R\,\ndof\bigr)
    \;=\; \cO\!\bigl(R\,\Ncell\,\nprim\bigr).
    \label{eq:fc3_pcp_params}
\end{align}

\paragraph{Self-energy kernel.} Identical in scaling to CP and INDSCAL, $\cO(R^{2}\ndof)$ per
frequency, but with the additional property that the reconstruction is exactly $S_3$-symmetric
in $(p(\mu),j,k)$, so ASR-sensitive observables receive only the quadratic-in-error contribution
(\cref{sec:res_fc}).


% ----------------------------------------------------------------------
\subsubsection{Further approximations}
\label{sub:fc3_other}
% ----------------------------------------------------------------------
The five ans\"atze of
\cref{sub:fc3_matsvd,sub:fc3_tucker,sub:fc3_cp,sub:fc3_indscal,sub:fc3_symcp} do not exhaust the
useful choices. Further approximations appear in the broader literature and can either be
combined with them or replace them in specific regimes.

\paragraph{Cross / CUR / skeleton approximations.} Rather than fitting a compressed form to a
precomputed $\Phi$, one can build the approximation from a small number of sampled fibers of the
tensor and interpolate the rest:
\begin{equation}
  \tilde\Phi_{\mu j k}^{\mathrm{CUR}}
  = \sum_{p,q,r}\Phi_{\mu,J_q,K_r}\,
    \Phi_{I_p,J_q,K_r}^{-1}\,\Phi_{I_p,j,k},
\end{equation}
where $I,J,K$ are index subsets of size $\sim R$ selected to maximise the volume of the
intersection subtensor, and $\Phi_{I,J,K}$ is inverted in the appropriate pseudoinverse sense.
This bypasses the construction of the full FC3 altogether, which is normally the computational
bottleneck (finite-displacement DFT). Rank-$R$ tensor cross
\cite{oseledets_tt-cross_2010, savostyanov_quasioptimality_2014} accesses $\cO(R^{2}\ndof)$ entries.
Adaptive cross approximation \cite{bebendorf_approximation_2000} and CUR-style subset selection
\cite{goreinov_theory_1997, mahoney_cur_2009} give analogous matrix-level
guarantees, and fiber-sampling Tucker (FSTD) of \citet{caiafa_generalizing_2010} yields an explicit
Tucker form from $\cO(R^{N-1})$ fibers. The tensor-CURT extension of
\citet{song_relative_2019} provides the first relative-error bounds in this setting. For
FC3 the main challenge is that cross methods have no intrinsic notion of ASR or $S_3$ symmetry.
A hybrid that samples in the symmetric CP parameter space would combine the two sets of
guarantees.

\paragraph{Tensor train and hierarchical Tucker.} The tensor-train (TT) format
\cite{oseledets_tensor-train_2011} writes an $N$-index tensor as a contraction of $N$ three-index
cores with bond ranks $(R_1,\dots,R_{N-1})$. For third-order tensors this collapses to
\begin{equation}
  \tilde\Phi_{\mu j k}^{\mathrm{TT}}
  = \sum_{\alpha_1,\alpha_2} G_1(\mu,\alpha_1)\,
    G_2(\alpha_1,j,\alpha_2)\,G_3(\alpha_2,k),
\end{equation}
which is structurally a Tucker1 SVD of the $(\mu)|(j,k)$ unfolding followed by a Tucker1 SVD of
the $(\alpha_1 j)|(k)$ matricisation of the residual. It does not offer independent advantages
beyond \cref{sub:fc3_tucker}. TT and the hierarchical Tucker (HT) format of
\citet{hackbusch_tensor_2019} become first-class methods at order
$\ge 4$ and are the natural representation for FC4 (four-phonon) tensors, where CP and Tucker
lose efficiency. \citet{lee_solving_2026} recently used Matrix Product States (= TT) to solve the
Peierls--Boltzmann equation in silicon.

\paragraph{Block-term decomposition (BTD).} \citet{de_lathauwer_decompositions_2008, de_lathauwer_decompositions_2008-1}
generalise both CP and Tucker by writing the tensor as a sum of low-multilinear-rank Tucker blocks,
\begin{equation}
  \tilde\Phi_{\mu j k}^{\mathrm{BTD}}
  = \sum_{r=1}^{R}\bigl(\mathcal{G}_r\times_1 A_r\times_2 B_r\times_3 C_r\bigr)_{\mu j k},
\end{equation}
with each $\mathcal{G}_r$ of small core shape (e.g.\ the rank-$(L,L,1)$ BTD has
$\mathcal{G}_r$ of shape $L\times L\times 1$). CP is the rank-$(1,1,1)$ special case; Tucker is
the single-block special case. For FC3 this structurally expresses block patterns like
per-atom-species or per-shell contributions in a single decomposition, and inherits the
uniqueness and ALS algorithms of the parent families.

\paragraph{Constrained / symmetry-adapted decompositions (CANDELINC).} CANDELINC
\cite{carroll_candelinc_1980} is CP with factor matrices constrained to user-chosen
linear subspaces,
\begin{equation}
  A=\Phi_A\hat A,\quad B=\Phi_B\hat B,\quad C=\Phi_C\hat C,
\end{equation}
where $\Phi_A,\Phi_B,\Phi_C$ are columnwise-orthonormal basis matrices. It reduces to an ordinary
CP fit in the reduced coordinates $(\hat A,\hat B,\hat C)$ applied to the projected tensor
$\hat\Phi = \Phi\times_1\Phi_A^\top\times_2\Phi_B^\top\times_3\Phi_C^\top$. For FC3 this is the
natural framework to couple any of the five ans\"atze with a space-group-compatible basis:
choose $\Phi_A,\Phi_B,\Phi_C$ to span the ASR-respecting, symmetry-invariant subspace (symfc
\cite{seko_projector-based_2024}, Fu--Marianetti \cite{fu_group_2019}, CSLD
\cite{zhou_lattice_2014, zhou_compressive_2019}), then fit INDSCAL or symmetric CP on the reduced
tensor.

\paragraph{Sparse / structured CP.} Imposing sparsity or smoothness priors on the CP factor
vectors yields decompositions whose rank-$R$ components are localised on a few atoms or DOFs.
As an example, one can penalise the loss with
\begin{equation}
  L(A,B,C) = \tfrac12\|\Phi - ( A,B,C)\|_{F}^{2}
            + \lambda_1\sum_r\bigl(\|a_r\|_1+\|b_r\|_1+\|c_r\|_1\bigr),
\end{equation}
optionally replacing $\|\cdot\|_1$ by group-sparse or total-variation terms
\cite{papalexakis_parcube_2012, kim_sparse_2007}, and are solved by
ALS or AO-ADMM.

% ----------------------------------------------------------------------
